%
Et bildekk blir fylt med luft i en garasje som holder $20\,^\circ$C. Det absolutte trykket i dekket er da \SI{3.0e5}{\pascal}, og volumet av lufta i dekket er \SI{0.030}{\metre\cubed}. Så kjøres bilen ut i kulda, der det er $-15\,^\circ$C. Anta at volumet av dekket ikke endres.
\begin{enumerate}[a)]
    \item Hvor mange luftmolekyler er det i dekket?
    \item Hva blir det absolutte trykket i dekket når lufta har fått samme temperatur som omgivelsene?
    \item En dekktrykkmåler viser \emph{overtrykket}, altså forskjellen mellom trykket i dekket og lufttrykket utenfor, $p_0 = \SI{1.0e5}{\pascal}$. Hva viser måleren før og etter? Hvor mange prosent har overtrykket sunket?
\end{enumerate}
%
